\section{Definición del atlas genealógico y convenciones de evaluación metrológica}

Este atlas no es una tabla de coincidencias. Su n\'ucleo matem\'atico conserva
cinco campos inseparables; el estatuto l\'ogico y el criterio de refutaci\'on se enuncian despu\'es, sin ocupar el lugar de ninguno de los cinco:

\begin{equation}
\boxed{
\begin{gathered}
\text{antecedente}\ \big|\ \text{ecuaci\'on}\ \big|\
\ \text{invariante conservado}\\
\text{significado}\ \big|\ \text{valor terminal}
\end{gathered}}
\qquad
\begin{gathered}
[\text{estatuto};\\ \text{criterio de refutación}]
\end{gathered}.
\label{eq:atlas-six-fields}
\end{equation}


\section{Corolario HMT de confluencia metrol\'ogica}
\label{sec:metrological-confluence}

\begin{definition}[Qu\'intupla metrol\'ogica]
Sea \(x\) una magnitud obtenida mediante una realizaci\'on HMT--MD. Su
\emph{qu\'intupla
metrol\'ogica} es
\[
 \mathfrak C(x)=
 \bigl(A_x\mid E_x\mid I_x\mid S_x\mid V_x\bigr),
\]
donde \(A_x\) es el antecedente tipado, \(E_x\) la ecuaci\'on de
determinaci\'on, \(I_x\) el invariante conservado por la realizaci\'on,
\(S_x\) el significado
estructural y \(V_x\) el valor terminal. No se permite formar \(V_x\) antes de
haber fijado los otros cuatro componentes. Dos salidas con el mismo decimal
no son iguales si difieren sus antecedentes o sus invariantes.
\end{definition}

\begin{corollary}[Confluencia metrol\'ogica]
\label{cor:metrological-confluence}
Las salidas de vac\'io, acci\'on, informaci\'on, temperatura, gravedad,
criticidad, espectro, mezcla y cosmolog\'ia construidas en los
\cref{ch:vacuum,ch:metrology,ch:chaos,ch:spectral,ch:gauss,ch:gravity,ch:modular,ch:ew,ch:cosmo}
admiten las qu\'intuplas de la \cref{tab:metrological-confluence}. En cada
fila, el valor de la quinta columna es consecuencia de las cuatro primeras y
no un dato que las selecciona. Las filas comparten nodos geneal\'ogicos
---modo \(30\), \(\log3\), determinante \(D_A\), canales \(90/120\) y cambio
\(6\to8\)---, pero conservan separados sus dominios y operadores.
\end{corollary}

\begingroup
\setlength{\tabcolsep}{2pt}
\renewcommand{\arraystretch}{1.22}
\scriptsize
\begin{longtable}{@{}
>{\raggedright\arraybackslash}p{25mm}
>{\raggedright\arraybackslash}p{37mm}
>{\raggedright\arraybackslash}p{24mm}
>{\raggedright\arraybackslash}p{29mm}
>{\raggedright\arraybackslash}p{32mm}@{}}
\caption{Corolario de confluencia metrol\'ogica HMT: las cinco columnas son
una sola unidad probatoria.}\label{tab:metrological-confluence}\\
\toprule
Antecedente & Ecuaci\'on & Invariante conservado & Significado estructural &
Valor terminal\\
\midrule
\endfirsthead
\multicolumn{5}{l}{\small\itshape Continuaci\'on del corolario de
confluencia metrol\'ogica HMT}\\
\toprule
Antecedente & Ecuaci\'on & Invariante conservado & Significado estructural &
Valor terminal\\
\midrule
\endhead
\midrule
\multicolumn{5}{r}{\itshape Contin\'ua en la p\'agina siguiente}\\
\endfoot
\bottomrule
\endlastfoot

\multicolumn{5}{@{}l}{\bfseries I. Acci\'on}\\*
dos realizaciones de la recta de acci\'on
(\cref{eq:action-twoway,eq:action-rapidity}) &
\(R_{\rm act}=H_5/\eta_{\rm ret}\),
\(\xi_{\rm act}=\tfrac12\log R_{\rm act}\) &
D\'ecada y orientaci\'on de hoja &
La acci\'on absoluta y su cociente relativo corresponden a funcionales
distintos; la
d\'ecada no se usa como ajuste &
\(R_{\rm act}\), \(\xi_{\rm act}\)\\

\addlinespace
\multicolumn{5}{@{}l}{\bfseries II. Vac\'io electromagn\'etico}\\*
\(T=q_+P_++q_-P_-\), \(30\to90/120\)
(\cref{eq:vacuum-operator,eq:three-four-completion}) &
\(V=(I-T^{90})(I-T^{120})^{-1}\) &
Positividad, espectro de hojas y completaci\'on \(3\to4\) &
El vac\'io es un operador bicapa; \(90=3\cdot30\) y \(120=4\cdot30\)
determinan el cuarto sector sin introducir una constante SI &
\(r_+=0.9999996285\ldots\), \(r_-=0.9997241371\ldots\)\\

respuesta espectral
(\cref{eq:vacuum-four,eq:vacuum-traces}) &
\(\widehat\varepsilon=r_+^2\),
\(\widehat\mu=r_-^2\),
\(\widehat Z=r_-/r_+\),
\(\widehat c=(r_+r_-)^{-1}\) &
Producto y cociente de hojas, estables bajo \(V\mapsto V\otimes I_{9^m}\) &
Permisividad, permeabilidad, impedancia y propagaci\'on son cuatro funciones
del mismo operador positivo &
\(0.9999992571\ldots\), \(0.9994483505\ldots\),
\(0.9997245085\ldots\), \(1.0002763105\ldots\)\\

rectas constitutivas
(\cref{eq:vacuum-sections}) &
\(\mu\varepsilon=c^{-2}\), \(\mu/\varepsilon=Z^2\) &
Homogeneidad en \(L_S,L_C,L_Q\) &
La carta dimensional es una secci\'on posterior, no una unidad pegada a un
n\'umero adimensional &
\(Z,\mu,\varepsilon,c\) como secciones exactas\\

\((Z_{\rm vac},h_a,\alpha_{\HMT})\)
(\cref{eq:charge,eq:quantum-electrical-identities}) &
\(e_a^2=2\alpha h_a/Z\), \(R_KG_0=2\), \(K_J\Phi_0=1\) &
Producto bidireccional y orientaci\'on opuesta de flujo/frecuencia &
Los cuantos el\'ectricos descienden del vac\'io y la acci\'on; la carga SI no
act\'ua como selector &
\(e_a,R_K,G_0,\Phi_{0,a},K_{J,a}\)\\

\addlinespace
\multicolumn{5}{@{}l}{\bfseries III. Informaci\'on, temperatura y radiaci\'on}\\*
trit uniforme y estructura cronol\'ogica \(108\)
(\cref{eq:boltzmann-clock,eq:log3-triple}) &
\(k_B\Theta_{\rm clk}\log3=h_{\rm int}/(108t_0)=E_P\) &
Entrop\'ia de una decisi\'on ternaria y periodo periodicidad temporal de (108) fases &
Un mismo cociente mide informaci\'on, temperatura y energ\'ia sin identificar
sus tres espacios de estados &
\(\log3=1.0986122886681098\ldots\)\\

\(G\hbar\) y periodicidad temporal de (108) fases
(\cref{eq:planck-energy,eq:planck-trit}) &
\(E_P=\pi\hbar/(54t_0)=k_B\Theta_{\rm clk}\log3\) &
Producto bidireccional \(G\hbar\) y contenido informacional TRIT &
El cuanto geom\'etrico de Planck y el cuanto informacional asociado a un trit
confluyen en una misma secci\'on de energ\'ia &
\(E_P=h/(108t_0)\)\\

espectro de Planck
(\cref{eq:wien-lambda,eq:wien-frequency}) &
\(\lambda_{\max}/\ell_0=108\log3/x_5\),
\(\nu_{\max}t_0=x_3/(108\log3)\) &
Mismo espectro, jacobianos distintos &
Los m\'aximos por longitud y por frecuencia son realizaciones funcionales de un antecedente común, no dos
aproximaciones de la misma coordenada &
\(23.896756779\ldots\), \(0.0237794888\ldots\)\\

continuo bos\'onico
(\cref{eq:stefan-hmt,eq:photon-number,eq:photon-entropy}) &
\(n_\gamma\ell_P^3=2\zeta(3)/[\pi^2(\log3)^3]\) y leyes \(T^4,T^3\) &
Ocupaci\'on bos\'onica y contenido informacional por trit &
Radiancia, n\'umero, energ\'ia y entrop\'ia forman una familia t\'ermica; no
son cuatro calibraciones independientes &
Coeficientes exactos en \(\pi,\zeta(3),\log3\)\\

\addlinespace
\multicolumn{5}{@{}l}{\bfseries IV. Gravedad, \`area y memoria modular}\\*
cierre circular periodicidad temporal de (108) fases
(\cref{thm:gravity-circular-selector,eq:gravity-theta108}) &
\(G=\theta_{108}c^5t_0^2/\hbar\),
\(\theta_{108}=(54/\pi)^2\) &
Igualdad del radio Compton y el radio gravitatorio &
La constante gravitatoria es una secci\'on seleccionada por un cierre de
longitudes, no un n\'umero injertado &
\(G_{\rm int}=(54/\pi)^2c^5t_0^2/\hbar\)\\

\(A_5/C_5\to P_5\)
(\cref{thm:gravity-p5-coupling,eq:gravity-G5-spectrum}) &
\(\mathcal G_5=GP_5\), \(P_5^2=P_5\), \(\Tr P_5=5\) &
Rango central y reducci\'on TT &
Cinco componentes del portador tensorial preceden a las dos helicidades
f\'isicas; no se confunden con ellas &
\(\Spec\mathcal G_5=\{G^{(5)},0^{(7)}\}\)\\

cambio \(6\to8\), derivación de Barbero--Immirzi y canales \(90/120\)
(\cref{eq:modular-area-operator,eq:modular-68-two-readers}) &
\(A_{\rm phys}/\ell_P^2=\gammaH a_0(N_{90}+N_{120})\) &
Cuanto de \`area, orientaci\'on hexada--octada y \(\ell_P^2\) &
\(\gammaH\) fija el cuanto geom\'etrico; esta monograf\'ia adopta dicha
magnitud sin reabrir su derivaci\'on &
\(\gammaH a_0=0.0016755940749\ldots\)\\

ocupaci\'on \(N\) y memoria \(D\)
(\cref{thm:modular-area-memory,eq:modular-inverse-channels}) &
\(D=90N_{90}+120N_{120}\),
\(N_{90}=4N-D/30\), \(N_{120}=D/30-3N\) &
Determinante \(30\) y procedencia de canal &
El \`area cuantifica la ocupaci\'on total; el Hamiltoniano hologr\'afico
modular conserva el sector \(90/120\) de procedencia &
\((N_{90},N_{120})\) reconstruido exactamente\\

estado de borde
(\cref{thm:modular-tfd,eq:modular-oriented-information,thm:modular-first-law}) &
\(I_{\rm or}=36\log3-S(\rho_A)\),
\(\delta S=\beta_{90}\delta\langle\mathcal A_{90}\rangle+
\beta_{120}\delta\langle\mathcal A_{120}\rangle\) &
Espectro de Schmidt y raz\'on \(\beta_{120}/\beta_{90}=4/3\) &
Entrop\'ia de entrelazamiento, informaci\'on orientada y \`area son tres
lecturas compatibles del mismo borde &
\(I_{\rm or}=0.0016691832\ldots\),
\((\beta_{90},\beta_{120})=(5.97265\ldots,7.96353\ldots)\)\\

red local infinita
(\cref{thm:modular-typeIII,eq:modular-typeIII-lambda}) &
\(\lambda_\partial=\exp(-30\Dmod)\) &
Grupo de razones de pesos bajo persistencia \(90/120\) &
El l\'imite es un generador modular GNS, no una matriz densidad global
traza--clase en la UHF &
\(\lambda_\partial=0.996669646468957\ldots\)\\

\addlinespace
\multicolumn{5}{@{}l}{\bfseries V. Criticidad y modo \(30\)}\\*
fases \(3m-1\)
(\cref{eq:pi-cancellation,eq:chaos-closed,eq:899}) &
\(u_0=1-\cos(37y)\sin(36y)/(12\sin3y)\),
\(899=30^2-1\) &
Normalizaci\'on cuadr\'atica y cancelaci\'on de \(\pi\) antes de iterar &
La constante de criticidad se deriva de un perfil cerrado cuyo modo \(30\)
separa asimismo \(90\) de \(120\) &
\(u_0(x)=x^2-(715769/2424603)x^4+O(x^6)\)\\

sucesi\'on superestable
(\cref{eq:negative-schwarzian,eq:renormalization}) &
\(\mathcal R(f)(x)=-a^{-1}f(f(-ax))\) &
Unimodalidad, Schwarziano negativo y combinatoria de duplicaci\'on &
Los cocientes finitos certifican la entrada en la clase universal; la
hiperbolicidad del punto fijo sigue siendo una obligaci\'on separada &
\(\delta_8=4.66921218915\ldots\),
\(\alpha_8=2.50296847988\ldots\)\\

\addlinespace
\multicolumn{5}{@{}l}{\bfseries VI. Espectro, aritm\'etica y cilindros}\\*
cuarto de torsi\'on
(\cref{eq:J-square,eq:catalan}) &
\(J^2=-I\), \(G=L(2,\chi_{-4})\) &
Orientaci\'on compleja y multiplicatividad del car\'acter &
Catalan es el segundo momento del operador de cuarto de torsi\'on, no un
decimal asociado por semejanza &
\(G=0.915965594177219\ldots\)\\

producto CRT
(\cref{eq:chi12,eq:L1chi12,eq:euler-kronecker}) &
\(\chi_{12}=\chi_{-3}\chi_{-4}\),
\(\gamma_{\mathbb Q(\sqrt3)}=\gamma+L'/L(1,\chi_{12})\) &
Orientaci\'on tri\'adica, cuarto de torsi\'on y t\'ermino finito de Dedekind &
La ret\'icula \(3/4/12\) determina un campo real y su constante de borde &
\(L(1)=0.7603459963\ldots\),
\(\gamma_{\mathbb Q(\sqrt3)}=1.05396560825\ldots\)\\

zeta tri\'adica renormalizada
(\cref{eq:bendersky,eq:odd-zeta-bendersky}) &
\(A_k=\exp(H_kB_{k+1}/(k+1)-\mathcal Z_3'(-k))\) &
Nivel de derivada zeta y normalizaci\'on Bernoulli &
Bendersky--Glaisher organiza \(\zeta(2n+1)\) como una familia graduada de
determinaciones &
\(A_0=\sqrt{2\pi}\), \(A_1=A_{\rm GK}\), familia \(A_k\)\\

cilindros cofinales
(\cref{eq:khinchin,eq:levy,eq:gauss-entropy,eq:lochs}) &
\(L=\pi^2/(12\log2)\), \(h_G=2L\),
\(\Lambda_b=\log b/h_G\) &
Medida de Gauss, crecimiento de denominadores y ritmo de codificaci\'on &
Khinchin, L\'evy y Lochs son observables distintos del mismo sistema
transportado &
\(K=2.6854520011\ldots\), \(L=1.1865691104\ldots\),
\(\Lambda_{10}=0.9702701144\ldots\)\\

operador de transferencia de Gauss
(\cref{eq:gkw-operator}) &
\(\mathcal L_Gf(x)=\sum_{n\ge1}(n+x)^{-2}
f((n+x)^{-1})\) &
Masa, medida invariante y complemento de media cero &
El autovalor subdominante mide la relajaci\'on de los cilindros; no se
selecciona desde un decimal tabulado &
\(\lambda_{\rm GKW}\): intervalo certificado por construir\\

ciclos irreducibles y vacancias
(\cref{eq:meissel-mertens,eq:twin-prime-constant,eq:vacancy-gamma,eq:vacancy-stieltjes}) &
\(B_1=\lim(\sum_{p\le x}p^{-1}-\log\log x)\),
\(\gamma_{\rm vac,j}^+=\varepsilon_{\rm vac}\gamma_j+\cdots\) &
Factorizaci\'on prima, discrepancia acotada y holonom\'ia de vacancia &
Suma local, densidad multiplicativa y t\'erminos finitos conservan
construcciones antecedentes distintas &
\(B_1=0.2614972128\ldots\), \(C_2=0.6601618158\ldots\),
\(\varepsilon_{\rm vac}=0.04575749056\ldots\)\\

levantamiento \(r^2=2\pmod{3^k}\)
(\cref{eq:pell-order}) &
\(u=3+2r\), \(\operatorname{ord}(u)=4\cdot3^{k-1}\) &
Operador de cuarto de torsi\'on compatible a toda profundidad non\'adica &
La unidad de Pell conserva simult\'aneamente orientaci\'on finita y memoria
profinita &
\(3+2\sqrt2=5.82842712474619\ldots\)\\

\addlinespace
\multicolumn{5}{@{}l}{\bfseries VII. Cambio de hoja, mezcla y electrod\'ebil}\\*
\(C^*\mapsto-C^*\), \(6\to8\)
(\cref{eq:modular-sheet-parities,eq:modular-ckm-coordinates}) &
\((A,C^*)\mapsto(A,-C^*)\),
\(-(8-6)/3=-2/3\) &
Paridad de hoja y defecto orientado hexada--octada &
Barbero y CKM son realizaciones distintas del cambio de estado, no
definiciones retrospectivas mutuas &
\(A,C^*\) recuperados racionalmente de los tres \`angulos\\

matriz de coordenadas \(M\)
(\cref{thm:modular-ckm-reflection,eq:modular-Rckm-matrix}) &
\(R_{\rm CKM}=M\operatorname{diag}(1,-1,1)M^{-1}\) &
Involuci\'on racional, \(R^2=I\), \(\det R=-1\) &
El espacio de mezcla conserva exactamente la coordenada que experiment\'o
la inversi\'on de hoja;
la fase CP queda fija en esta involuci\'on &
Matriz racional exacta, \(\det M=-3857/6\)\\

operador angular \(T\)
(\cref{thm:ew-weinberg,eq:ew-sin2thetaW}) &
\(D_A=\det T=e^{-2A_{\rm rad}}\),
\(\sin^2\theta_W=1-D_A\) &
Determinante par y adyacencia de rutas \(W/Z\) &
El vac\'io y Weinberg comparten operador angular, pero determinan funciones
espectrales diferentes &
\(D_A=0.775129119247\ldots\),
\(\sin^2\theta_W=0.224870880753\ldots\)\\

carta de calibre en aproximaci\'on de \`arbol
(\cref{eq:ew-g,eq:ew-gprime,eq:ew-GF,cor:ew-Higgs-lambda}) &
\(g=2\sqrt{\pi\alpha/(1-D_A)}\),
\(G_F^{(0)}=\pi\alpha/[\sqrt2(1-D_A)m_W^2]\) &
El mismo \(D_A\) conserva la separaci\'on de canal neutro/cargado &
Los acoplamientos, la constante de Fermi y el autoacoplamiento de Higgs
forman una composici\'on tipada; las correcciones radiativas no se incorporan
al funcional interno &
\(G_F^{(0)}=1.1157162817\cdot10^{-5}\,\mathrm{GeV}^{-2}\),
\(\lambda_H=0.123847911548\ldots\)\\

funcionales \((\alpha,\varphi,1000,55)\)
(\cref{eq:atlas-anom-mu}) &
\(a_\mu=\alpha/(2\pi)+\alpha(\varphi-1)/1000+\alpha^2/(55000)\) &
Suma de tres contribuciones con coeficientes congelados &
Control de composici\'on lept\'onica; no sustituye la serie radiativa completa
de QED &
\(a_\mu=0.001165920713\ldots\)\\

\addlinespace
\multicolumn{5}{@{}l}{\bfseries VIII. Cosmolog\'ia y controles de alcance}\\*
radio \(R_{108}=\ell_P\)
(\cref{eq:cosmo-ds-H-Lambda,eq:cosmo-ds-temperature,eq:cosmo-ds-entropy}) &
\(H_{108}=\pi/(54t_0)\),
\(\Lambda_{108}=\pi^2/(972\ell_0^2)\) &
Un solo radio bajo la realizaci\'on de Sitter &
Expansi\'on, curvatura, temperatura y entrop\'ia son cuatro magnitudes de
una misma geometr\'ia condicionada &
\(T_{\rm dS}/\Theta_{\rm clk}=\log3/(2\pi)\),
\(S_{\rm dS}/k_B=\pi\)\\

escala de Perron
(\cref{eq:cosmo-perron-a,eq:cosmo-minus-six,eq:cosmo-EC-scaling}) &
\(a_n/a_0=\varphi^{n/3}\),
\(M_n=(a_n/a_0)^{-6}\) &
Modo estable \(\varphi^{-2n}\) y diluci\'on cuadr\'atica de esp\'in &
El exponente \(-6\) representa la memoria activa en la ley cosmol\'ogica; la
amplitud Einstein--Cartan requiere su transductor &
\(M_na_n^6=a_0^6\),
\(H_\varphi t_0=0.1604039416865\ldots\)\\

controles de composici\'on
(\cref{eq:atlas-bohr,eq:atlas-born}) &
\(a_0=\hbar/(m_ec\alpha)\),
\(E_*=h_{\rm int}c_{\rm int}/(e_{\rm int}\ell_0^2)\) &
Homogeneidad dimensional y no alimentaci\'on por nombre &
Bohr se compone a partir de magnitudes definidas fuera del alcance de esta
monograf\'ia; Born--Infeld exige a\'un
una acci\'on determinantal que seleccione su escala &
Secciones exactas \(a_0,E_*\), sin promover nuevo selector\\

\end{longtable}
\endgroup

\begin{proof}
Cada fila se obtiene sustituyendo exclusivamente objetos ya construidos en la
ecuaci\'on citada en su primera o segunda columna. En vac\'io, el c\'alculo
es espectral: el teorema funcional para \(V\) produce \(r_\pm\), y las cuatro
constitutivas son funciones de esos dos autovalores. En la realizaci\'on
t\'ermica, la estructura cronol\'ogica \(108\) y la entrop\'ia uniforme de
un trit producen el producto
\(k_B\Theta_{\rm clk}\), tras lo cual Planck, Wien y radiaci\'on son
composiciones. En gravedad, la condici\'on circular fija
\(\theta_{108}\), mientras \(P_5\) y la reducci\'on TT fijan el tipo del
portador. En la realizaci\'on modular, el sistema lineal de determinante \(30\)
reconstruye ambos canales y la descomposici\'on espectral del estado reducido
produce entrop\'ia y primera ley. Las construcciones de criticidad,
caracteres, Gauss y
Pell se prueban en sus respectivos operadores. La reflexi\'on CKM se obtiene
por conjugaci\'on de una involuci\'on diagonal, y el sector electrod\'ebil por
la funci\'on determinante del operador angular. Finalmente, las dos
realizaciones cosmol\'ogicas emplean, respectivamente, un radio periodicidad temporal de (108) fases y el
autovalor estable
de Perron. En ning\'un paso interviene el decimal de la quinta columna para
elegir el antecedente. Esto prueba la terminalidad y la confluencia declaradas.
\end{proof}

\section{Relaciones estructurales principales de confluencia geneal\'ogica}

El corolario anterior no se limita a enumerar filas: determina qu\'e
operadores comparten un antecedente y el punto exacto en que sus dominios se
separan. Las cinco configuraciones siguientes contienen las relaciones
causales de mayor densidad del volumen.

\subsection{Modo \(30\): completaci\'on arm\'onica, criticidad y memoria}

El mismo entero aparece en tres construcciones exactas:
\begin{equation}
\begin{aligned}
90&=3\cdot30, &120&=4\cdot30,
&\frac{1-q^{90}}{1-q^{120}}
 &=\frac{1+s+s^2}{1+s+s^2+s^3},\qquad s=q^{30},\\
899&=30^2-1=29\cdot31,
&&&D&=90N_{90}+120N_{120}.
\end{aligned}
\label{eq:atlas-mode30-confluence}
\end{equation}
La primera l\'inea es la completaci\'on \(3\to4\) del operador de vac\'io; la
segunda normaliza el germen dodecaf\'asico sin introducir \(\pi\) en la
iteraci\'on; la tercera conserva cu\'antas ocupaciones llegaron por cada canal
modular. El invariante com\'un es la partici\'on del modo \(30\); los
codominios son diferentes. Por ello, el cruce es estructural y no una
identificaci\'on de \(V\), \(u_0\) y \(\HHM\).

\paragraph{Identidades de completación del modo \(30\), invertibilidad y condición de incompatibilidad.}
Las tres identidades de \cref{eq:atlas-mode30-confluence} son exactas. Queda
invalidada la confluencia si \(90/120\) dejan de ser m\'ultiplos consecutivos del
mismo modo, si \(899\ne30^2-1\), o si la inversi\'on
\((N,D)\mapsto(N_{90},N_{120})\) pierde su determinante \(30\). No la falsan
las diferencias de espectro: esas diferencias son precisamente las que
mantienen tipadas las tres realizaciones.

\subsection{\(\log3\): invariante residual, informaci\'on y temperatura}

La confluencia informacional se resume en la cadena
\begin{equation}
S_{\TRIT}=k_B\log3,\qquad
E_P=k_B\Theta_{\rm clk}\log3,\qquad
I_{\rm or}=36\log3-S(\rho_A).
\label{eq:atlas-log3-confluence}
\end{equation}
El primer miembro mide una decisi\'on ternaria uniforme; el segundo determina
su magnitud energ\'etica en la estructura cronol\'ogica \(108\); el tercero compara la
capacidad de \(36\) trits con la entrop\'ia del estado reducido de borde.
El escalar se conserva,
pero la unidad, la temperatura y el estado reducido permanecen tipol\'ogicamente
separados. Esta
separaci\'on es lo que permite leer de forma precisa la identidad
\(E_P=k_B\Theta_{\rm clk}\log3\): la realizaci\'on de Planck aporta el cuanto
de energ\'ia, la realizaci\'on TRIT su contenido informacional y la estructura
cronol\'ogica el transductor.

\paragraph{Conservación de \(\log 3\), separación dimensional y condición de incompatibilidad.}
La triple lectura es una composici\'on exacta de
\cref{eq:log3-triple,eq:planck-trit,eq:modular-oriented-information}. Se
invalida si se reemplaza \(\log3\) por una entrop\'ia binaria, si
\(\Theta_{\rm clk}\) se fija desde kelvin antes de la identidad, o si
\(S(\rho_A)\) se calcula con un estado distinto del obtenido por traza parcial.

\subsection{\texorpdfstring{\(D_A\)}{Determinante angular}: determinante de hojas y relaci\'on angular de Weinberg}

El operador angular \(T=q_+P_++q_-P_-\) posee dos familias funcionales:
\begin{equation}
F(T)^2\longrightarrow
(\widehat\varepsilon,\widehat\mu,\widehat Z,\widehat c),
\qquad
\det T=q_+q_-=D_A\longrightarrow
\sin^2\theta_W=1-D_A.
\label{eq:atlas-DA-confluence}
\end{equation}
La primera conserva el espectro completo de la respuesta de vac\'io; la
segunda descarta la orientaci\'on individual de las hojas y conserva s\'olo su
producto. De ah\'i que \(D_A\) sea par bajo \(C^*\mapsto-C^*\) y pueda
determinar el cociente \(W/Z\), mientras la impedancia conserva el cociente
de autovalores. La relaci\'on contiene m\'as informaci\'on que una
coincidencia angular: es
una bifurcaci\'on funcional exacta del mismo operador antecedente.

\paragraph{Exactitud del determinante angular y obstrucciones de la interfaz electrodébil.}
El lado espectral y la identidad \(D_A=e^{-2A_{\rm rad}}\) son exactos en la
carta declarada. La realizaci\'on de calibre conserva su estatuto de interfaz hasta
construir el entrelazador de bases. Un cociente de rutas no adyacentes, una
dependencia impar de \(C^*\) o \((m_W/m_Z)^2\ne D_A\) falsan esa interfaz sin
afectar al operador de vac\'io.

\subsection{Sectores \(90\) y \(120\): vac\'io, \`area y din\'amica modular}

Los canales \(90\) y \(120\) tienen tres funciones no intercambiables:
son exponentes de propagaci\'on en \(V\), multiplicidades de la incidencia
\(6\to8\), y pesos espectrales en \(\HHM\). Su raz\'on \(4/3\) reaparece en
las temperaturas modulares,
\[
\frac{120}{90}=\frac43=\frac{\beta_{120}}{\beta_{90}},
\]
pero el operador de vac\'io no es el Hamiltoniano modular y una temperatura
modular no es una permeabilidad. La confluencia consiste en que la misma
genealog\'ia conserva tres magnitudes: propagaci\'on, ocupaci\'on areal y
memoria de canal.

\paragraph{Reconstrucción finita de los canales \(90/120\) e hipótesis del límite modular.}
La reconstrucci\'on finita y la raz\'on \(4/3\) son exactas. La clasificaci\'on
de tipo III requiere persistencia infinita de ambos canales. Desaparecer uno
de ellos a partir de alg\'un nivel, perder compatibilidad de estados o obtener
un grupo de razones distinto de
\(\{\lambda_\partial^n:n\in\Z\}\) falsar\'ia esa clasificaci\'on, no la
inversi\'on finita.

\subsection{Par\'ametro de Barbero--Immirzi, entrop\'ia y CKM como realizaciones del tr\'ansito \(6\to8\)}

La interfaz de confluencia completa se escribe sin convertir una
realizaci\'on en causa de las restantes:
\begin{equation}
(6\to8,\,90/120,\,C^*\mapsto-C^*)\longrightarrow
\begin{cases}
\gammaH\ell_P^2\longrightarrow A_{\rm phys},\\
\rho_A\longrightarrow(-\log\rho_A,S,I_{\rm or},\delta S),\\
C^*\mapsto-C^*\longrightarrow
R_{\rm CKM}=M\operatorname{diag}(1,-1,1)M^{-1}.
\end{cases}
\label{eq:atlas-bic-entropy-ckm}
\end{equation}
La primera realizaci\'on fija el cuanto geom\'etrico; la segunda mide la
informaci\'on del estado reducido y su respuesta areal; la tercera registra la
orientaci\'on de hoja en
el espacio de mezcla. En paralelo,
\(E_P=k_B\Theta_{\rm clk}\log3\) fija el cuanto informacional de energ\'ia.
As\'i gravedad, entrop\'ia e informaci\'on se ecuacionan mediante
\((\ell_P^2,E_P)\), sin identificar \(\gammaH\), \(k_B\) o los \`angulos CKM.

\paragraph{Independencia causal de Barbero--Immirzi, información modular y mezcla CKM.}
La salida \(\gammaH\) es un resultado antecedente demostrado de su derivación correspondiente; el
operador de \`area, la purificaci\'on, la primera ley local y la reflexi\'on CKM
son exactos en esta monograf\'ia. La clasificaci\'on modular infinita y la lectura
f\'isica de horizonte conservan sus hip\'otesis. El diagrama queda invalidado si
se usa CKM para recalibrar \(\gammaH\), si se iguala una hexada/octada de
incidencia con caras/v\'ertices sin entrelazador, o si
\(R_{\rm CKM}^2\ne I\) o \(\det R_{\rm CKM}\ne-1\).

\section{Independencia teoremática del corolario de confluencia metrológica}

La \cref{tab:metrological-confluence} no remite a una tabla exterior para
completar ninguna de sus cinco columnas. Cada antecedente se define en el
cuerpo; cada ecuaci\'on est\'a demostrada o tipada en el localizador citado;
cada invariante declara qu\'e se conserva y tambi\'en qu\'e se descarta; cada
significado impide la identificaci\'on por decimal; y cada valor terminal se
obtiene despu\'es de esa cadena. Los estatutos y criterios de refutación permanecen en las
fichas extensas que siguen, donde se desarrollan sin la compresi\'on inevitable
de una tabla sin\'optica.

En consecuencia, la tabla no sustituye el desarrollo: lo complementa. Puede leerse
de forma sinóptica para reconocer la confluencia, pero su validez procede de los
teoremas y pruebas del cuerpo, y cada fila dispone a continuaci\'on de una
ficha aut\'onoma con antecedente, ecuaci\'on, significado, valor, estatuto y
control negativo.

\section{Acci\'on, vac\'io y constantes cu\'antico--el\'ectricas}

\subsection*{Secciones bidireccionales de la acci\'on }
\textbf{Antecedente.} Dos realizaciones orientadas de una misma recta de
acci\'on. \textbf{Ecuaci\'on.}
\[
R_{\rm act}=\frac{\hbar_{\rm pre}}{\hbar_{\rm ret}}
=\frac{H_5}{\eta_{\rm ret}},
\qquad \xi_{\rm act}=\tfrac12\log R_{\rm act}.
\]
\textbf{Estructura y significado.} La d\'ecada com\'un fija el tipo; el cociente
conserva la orientaci\'on. \textbf{Valor.} \(R_{\rm act}\) y
\(\xi_{\rm act}\) en la carta interna, sin comparar mantisas dimensionales con
residuos. \textbf{Estatuto.} \emph{Teorema interno exacto}.
\textbf{Criterio de refutación.} Usar la d\'ecada com\'un como corrector relativo o mezclar
las dos secciones.

\subsection*{Operador constitutivo del vacío: espectro positivo y reconstrucción conjunta de \(\varepsilon,\mu,Z,c\)}
\textbf{Antecedente.} El operador bicapa \(T=q_+P_++q_-P_-\) y la genealog\'ia
\(30\to90/120\). \textbf{Ecuaci\'on.}
\[
V=(I-T^{90})(I-T^{120})^{-1},
\qquad
\frac{1-q^{90}}{1-q^{120}}
=\frac{1+s+s^2}{1+s+s^2+s^3},\quad s=q^{30}.
\]
Si \(\Spec V=\{r_+,r_-\}\),
\[
\widehat\varepsilon=r_+^2,
\quad\widehat\mu=r_-^2,
\quad\widehat Z=r_-/r_+,
\quad\widehat c=(r_+r_-)^{-1}.
\]
\textbf{Estructura y significado.} Las cuatro constitutivas son funciones de
un solo operador positivo; cada hoja completa tres tramos por un cuarto tramo.
La amplificaci\'on \(V_m=V\otimes I_{9^m}\) conserva los funcionales bajo
la expectativa tracial.
\textbf{Valor.}
\begin{align*}
\widehat\varepsilon&=0.9999992570966367027604416155\ldots,\quad
\widehat\mu=0.9994483504793269350899815559\ldots,\\
\widehat Z&=0.9997245085389372154555543466\ldots,\quad
\widehat c=1.0002763104861575261063862689\ldots.
\end{align*}
\textbf{Estatuto.} \emph{Teorema interno exacto}; amplificaci\'on non\'adica:
\emph{Formalizaci\'on nueva exacta}. \textbf{Criterio de refutación.} P\'erdida de
positividad, residuo en la identidad \(3\to4\), o variaci\'on de los
invariantes bajo la expectativa tracial.

\subsection*{Secciones constitutivas de \(\varepsilon\), \(\mu\), \(Z\) y \(c\) en las cartas orientadas}
\textbf{Antecedente.} Los cuatro coeficientes anteriores y el marco
\((u_S,u_C,u_T,u_Q)\). \textbf{Ecuaci\'on.}
\[
Z=\widehat Zu_Su_Q^{-2},\quad c=\widehat cu_C,\quad
\mu=\widehat\mu u_Su_C^{-1}u_Q^{-2},\quad
\varepsilon=\widehat\varepsilon u_S^{-1}u_C^{-1}u_Q^2.
\]
\textbf{Estructura y significado.} La permisividad y la permeabilidad no
reciben unidades despu\'es de un c\'alculo adimensional: son secciones de rectas
constitutivas. \textbf{Valor.} Las cuatro secciones anteriores, con
\(\mu\varepsilon=c^{-2}\) y \(\mu/\varepsilon=Z^2\).
\textbf{Estatuto.} \emph{Teorema interno exacto}. \textbf{Criterio de refutación.} Introducir
\(c,\mu_0,\varepsilon_0\) SI antes del operador o violar la homogeneidad.

\subsection*{Derivación conjunta de la carga, la resistencia de von Klitzing, la conductancia, el flujo y la constante de Josephson}
\textbf{Antecedente.} \((Z_{\rm vac},h_a,\alpha_{\HMT})\).
\textbf{Ecuaci\'on.}
\[
e_a=\sqrt{\frac{2\alpha_{\HMT}h_a}{Z_{\rm vac}}},\quad
R_K=\frac{Z_{\rm vac}}{2\alpha_{\HMT}},\quad
G_0=\frac{4\alpha_{\HMT}}{Z_{\rm vac}},
\]
\[
\Phi_{0,a}=\sqrt{\frac{h_aZ_{\rm vac}}{8\alpha_{\HMT}}},
\qquad K_{J,a}=\Phi_{0,a}^{-1}.
\]
\textbf{Estructura y significado.} \(R_K\) y \(G_0\) son invariantes
bidireccionales; \(\Phi_0\) y \(K_J\) conservan la orientaci\'on de acci\'on con
exponentes opuestos. \textbf{Valor.} Las expresiones exactas en sus rectas;
\(R_KG_0=2\) y \(K_J\Phi_0=1\). \textbf{Estatuto.}
\emph{Teorema interno exacto}. \textbf{Criterio de refutación.} Usar la carga SI como selector o
fallar \(\alpha=Ze^2/(2h)\).

\section{Constantes t\'ermicas y radiaci\'on}

\subsection*{Relaci\'on de Boltzmann y contenido entr\'opico TRIT }
\textbf{Antecedente.} Decisi\'on ternaria uniforme, secci\'on de acci\'on y
estructura cronol\'ogica \(108\). \textbf{Ecuaci\'on.}
\[
k_B\Theta_{\rm clk}\log3=\frac{h_{\rm int}}{108t_0}
=\frac{2\pi\hbar_{\rm int}}{108t_0},
\qquad
\log3=\frac{S_{\TRIT}}{k_B}=\frac{E_P}{k_B\Theta_{\rm clk}}.
\]
\textbf{Estructura y significado.} \(\log3\) es simult\'aneamente defecto
espectral, contenido de una decisi\'on y raz\'on t\'ermica; los operadores
permanecen
distintos. \textbf{Valor.} \(\log3=1.0986122886681098\ldots\) y la secci\'on
energ\'ia--temperatura anterior. \textbf{Estatuto.} \emph{Confluencia
exacta}. \textbf{Criterio de refutación.} Separar arbitrariamente \(k_B\) de
\(\Theta_{\rm clk}\) o inferir identidad de dominios por igualdad del escalar.

\subsection*{Ley de desplazamiento de Wien }
\textbf{Antecedente.} El espectro de Planck evaluado en
\(\Theta_{\rm clk}\). \textbf{Ecuaci\'on.}
\[
x_5=5+W_0(-5e^{-5}),\quad
\frac{\lambda_{\max}}{\ell_0}=\frac{108\log3}{x_5},
\]
\[
x_3=3+W_0(-3e^{-3}),\quad
\nu_{\max}t_0=\frac{x_3}{108\log3}.
\]
\textbf{Estructura y significado.} Longitud de onda y frecuencia son dos
parametrizaciones espectrales con jacobianos distintos. \textbf{Valor.}
\(\lambda_{\max}/\ell_0=23.896756779\ldots\) y
\(\nu_{\max}t_0=0.0237794888\ldots\). \textbf{Estatuto.}
\emph{Teorema interno exacto}. \textbf{Criterio de refutación.} Confundir \(x_5\) con \(x_3\) o
introducir kelvin antes del producto t\'ermico.

\subsection*{Ley de Stefan--Boltzmann y gas de fotones }
\textbf{Antecedente.} La misma escala t\'ermica, el continuo bos\'onico y el
funcional \(\zeta(3)\). \textbf{Ecuaci\'on.}
\[
\mathcal F(\Theta_{\rm clk})
=\frac{2\pi^5}{15\,108^4(\log3)^4}
\frac{h_{\rm int}}{\ell_0^2t_0^2},
\]
\[
n_\gamma\ell_P^3=\frac{2\zeta(3)}{\pi^2(\log3)^3},
\qquad
\frac{u_\gamma\ell_P^3}{E_P}=\frac{\pi^2}{15(\log3)^4}.
\]
\textbf{Estructura y significado.} \(\zeta(3)\) cuenta ocupaci\'on bos\'onica;
\(\log3\) fija el contenido informacional TRIT por celda. \textbf{Valor.}
Los coeficientes
exactos anteriores. \textbf{Estatuto.} \emph{Composici\'on exacta con
estructura de partida expl\'icita}. \textbf{Criterio de refutación.} P\'erdida de factores
\(2\pi\), dimensiones de flujo incorrectas o fusi\'on de las procedencias matemáticas de
\(\zeta(3)\) y \(\log3\).

\section{Criticidad y constantes din\'amicas}

\subsection*{Perfil dodecaf\'asico y normalizaci\'on por \(899\)}
\noindent\textit{Resultados rectores: , y
.}\par
\textbf{Antecedente.} Doce fases \(3m-1\) y modo elemental \(30\).
\textbf{Ecuaci\'on.} Con \(y=x/\sqrt{899}\),
\[
u_0(x)=1-\frac1{12}\sum_{m=1}^{12}
\cos\!\bigl(2(3m-1)y\bigr)
=1-\frac{\cos37y\sin36y}{12\sin3y},
\]
\[
899=30^2-1=29\cdot31,\qquad
u_0(x)=x^2-\frac{715769}{2424603}x^4+O(x^6).
\]
\textbf{Estructura y significado.} La fase \(\pi\) se cancela antes de
iterar; el perfil resultante es anal\'itico, unimodal cuadr\'atico y de
Schwarziano negativo en \([-1,1]\). \textbf{Valor.} Perfil y coeficientes
racionales exactos. \textbf{Estatuto.} \emph{Teorema interno exacto}.
\textbf{Criterio de refutación.} Un segundo punto cr\'itico, Schwarziano no negativo o
residuo en la identidad trigonom\'etrica.

\subsection*{Aproximantes superestables de Feigenbaum y certificado hiperbólico de universalidad}
\textbf{Antecedente.} La sucesi\'on superestable del germen anterior.
\textbf{Ecuaci\'on.}
\[
\mathcal R(f)(x)=-a(f)^{-1}f\bigl(f(-a(f)x)\bigr),
\qquad a(f)=-f(1).
\]
\textbf{Estructura y significado.} Los cocientes finitos prueban que la
familia entra en la ruta adecuada; la universalidad requiere el punto fijo y
su hiperbolicidad. \textbf{Valor.}
\(\delta_8=4.66921218915\ldots\),
\(\alpha_8=2.50296847988\ldots\). \textbf{Estatuto.} Aproximantes:
\emph{Certificado finito}; l\'imite universal: \emph{Programa con
criterio de refutación}. \textbf{Criterio de refutación.} M\'as de un autovalor inestable, cola no acotada
o transversalidad nula.

\section{Constantes espectrales, aritm\'eticas y de vacancia}

\subsection*{Estructura compleja interna y constante de Catalan }
\textbf{Antecedente.} \(J^2=-I\) y el car\'acter
\(\chi_J(n)=i^{n-1}=\chi_{-4}(n)\) sobre impares. \textbf{Ecuaci\'on.}
\[
G=L(2,\chi_J)=\sum_{n\ge0}\frac{(-1)^n}{(2n+1)^2}.
\]
\textbf{Estructura y significado.} Catalan es el momento cuadr\'atico del
cuarto de torsi\'on. \textbf{Valor.} \(G=0.915965594177219\ldots\).
\textbf{Estatuto.} \emph{Teorema interno exacto}. \textbf{Criterio de refutación.} Falta de
multiplicatividad o discrepancia con \(\chi_{-4}\).

\subsection*{Car\'acter dodecaf\'asico y campo cuadr\'atico real }
\textbf{Antecedente.} Producto CRT
\(\chi_{12}=\chi_{-3}\chi_{-4}\). \textbf{Ecuaci\'on.}
\[
L(1,\chi_{12})=\frac{\log(2+\sqrt3)}{\sqrt3},
\qquad L(2,\chi_{12})=\frac{\pi^2}{6\sqrt3},
\]
\[
\gamma_{\mathbb Q(\sqrt3)}
=\gamma+\frac{L'(1,\chi_{12})}{L(1,\chi_{12})}.
\]
\textbf{Estructura y significado.} La orientaci\'on tri\'adica y el cuarto de
torsi\'on producen el car\'acter real de \(\mathbb Q(\sqrt3)\);
Euler--Kronecker es
el t\'ermino finito de su zeta de Dedekind. \textbf{Valor.}
\(0.7603459963009463\ldots\), \(0.9497031262940094\ldots\) y
\(1.05396560824848258\ldots\). \textbf{Estatuto.} \emph{Teorema interno exacto};
la identidad compuesta conserva las hipótesis y normalizaciones declaradas.
\textbf{Criterio de refutación.} Patr\'on de residuos distinto de \((+,-,-,+)\) o fallo de
la factorizaci\'on de Dedekind.

\subsection*{Familia Bendersky--Glaisher }
\textbf{Antecedente.} Funcional zeta tri\'adico renormalizado.
\textbf{Ecuaci\'on.}
\[
A_k=\exp\!\left(\frac{H_kB_{k+1}}{k+1}-\mathcal Z_3'(-k)\right),
\quad
\zeta(2n+1)=(-1)^{n+1}\frac{2(2\pi)^{2n}}{(2n)!}\log A_{2n}.
\]
\textbf{Estructura y significado.} Ap\'ery y los valores impares de zeta son
niveles de una torre, no coincidencias aisladas. \textbf{Valor.}
\(A_0=\sqrt{2\pi}\), \(A_1=A_{\rm GK}\), y la familia exacta anterior.
\textbf{Estatuto.} \emph{Teorema interno exacto}. \textbf{Criterio de refutación.} Fallo en
\(A_0,A_1\) o en la identidad de \(\zeta(2n+1)\).

\subsection*{Sistema dinámico de Gauss: \(K\), \(L\), \(\lambda_{\mathrm{GKW}}\) y cilindros transportados}
\textbf{Antecedente.} Evaluaci\'on de cilindros HMT y medida
\(d\mu_G=dx/[\log2(1+x)]\). \textbf{Ecuaci\'on.}
\[
\log K=\int\log a_1\,d\mu_G,\quad
L=-\int\log x\,d\mu_G=\frac{\pi^2}{12\log2},
\]
\[
h_G=2L=\frac{\pi^2}{6\log2},\qquad
\Lambda_b=\frac{6\log2\log b}{\pi^2}.
\]
\textbf{Estructura y significado.} Khinchin mide d\'igitos; L\'evy mide
crecimiento de denominadores; Lochs compara ritmos de codificaci\'on.
\textbf{Valor.}
\begin{align*}
K&=2.685452001065306\ldots,
&L&=1.18656911041562545\ldots,\\
h_G&=2.37313822083125091\ldots,
&\Lambda_{10}&=0.97027011439203393\ldots,\\
&&\Lambda_{729}&=2.77761896637430578\ldots.
\end{align*}
\textbf{Estatuto.} Cofinalidad e identidades: \emph{Exacto en el sistema
transportado}; constante GKW: \emph{Programa abierto con criterio de refutación}.
\textbf{Criterio de refutación.} No cofinalidad de cilindros, medida no invariante o ausencia
de gap espectral en el espacio declarado.

\subsection*{Ciclos aritméticos de Meissel--Mertens y series regularizadas de primos}
\textbf{Antecedente.} Selector de ciclos irreducibles y factorizaci\'on
exacta. \textbf{Ecuaci\'on.}
\[
B_1=\lim_{x\to\infty}\left(\sum_{p\le x}\frac1p-\log\log x\right),
\qquad
C_2=\prod_{p>2}\frac{p(p-2)}{(p-1)^2}.
\]
\textbf{Estructura y significado.} \(B_1\) es una suma local regularizada;
\(C_2\) una densidad multiplicativa de pares. \textbf{Valor.}
\(B_1=0.2614972128476428\ldots\),
\(C_2=0.6601618158468696\ldots\). \textbf{Estatuto.}
\emph{Demostrado con estructura de partida expl\'icita}: selector primo y
cota de cola declarados. \textbf{Criterio de refutación.} Usar ceros o decimales para elegir
los primos, o afirmar el l\'imite sin ejecutar la cota de cola.

\subsection*{Vacancias, holonom\'ia y constantes de Stieltjes }
\textbf{Antecedente.} Palabra mec\'anica de vacancias y discrepancia acotada.
\textbf{Ecuaci\'on.}
\[
\varepsilon_{\rm vac}=\log_{10}\frac{10}{9},\quad
\gamma_{\rm vac}^+=\lim_N\left(\sum_{n\le N}\frac{z_n}{n}
-\varepsilon_{\rm vac}\log N\right),
\]
\[
\rho_\varepsilon(k)=e^{2\pi ik\varepsilon_{\rm vac}},\qquad
\gamma_{\rm vac,j}^+=\varepsilon_{\rm vac}\gamma_j+
\sum_{n\ge1}\frac{(z_n-\varepsilon_{\rm vac})(\log n)^j}{n}.
\]
\textbf{Estructura y significado.} Densidad, t\'ermino finito y fase son tres
tipos; la torre prolonga el t\'ermino finito. \textbf{Valor.}
\(\varepsilon_{\rm vac}=0.0457574905606751\ldots\); los restantes se determinan
por sus l\'imites certificados. \textbf{Estatuto.} \emph{Teorema interno exacto con
discrepancia}. \textbf{Criterio de refutación.} P\'erdida de la cota de discrepancia o
discordancia del nivel \(j=0\).

\subsection*{Levantamiento non\'adico de Pell }
\textbf{Antecedente.} \(r^2=2\pmod{3^k}\). \textbf{Ecuaci\'on.}
\[
u=3+2r,\qquad \operatorname{ord}(u)=4\cdot3^{k-1},
\qquad \varprojlim\langle u\rangle=C_4\times\mathbb Z_3.
\]
\textbf{Estructura y significado.} El cuarto de giro persiste a toda
profundidad; la realizaci\'on arquimediana es \((1+\sqrt2)^2\).
\textbf{Valor.} \(3+2\sqrt2=5.82842712474619\ldots\).
\textbf{Estatuto.} \emph{Teorema interno exacto}. \textbf{Criterio de refutación.} Orden distinto
o levantamientos incompatibles.

\section{Gravedad, \'area e informaci\'on modular}

\subsection*{Constante gravitatoria y representaci\'on tensorial pentacomponente }
\textbf{Antecedente.} Estructura cronol\'ogica \(108\), acci\'on y proyector
central de rango
cinco. \textbf{Ecuaci\'on.}
\[
G_a(\theta)=\theta\frac{c_{\rm int}^5t_0^2}{\hbar_a},
\qquad \theta_{108}=\left(\frac{54}{\pi}\right)^2,
\qquad \mathcal G_5=G_a(\theta_{108})P_5.
\]
\textbf{Estructura y significado.} \(G\) es una secci\'on seleccionada por el
cierre Compton--gravedad; \(P_5\) determina un portador de cinco componentes,
que la reducci\'on TT lleva a dos helicidades sin masa.
\textbf{Valor.} \(G_{\rm int}=\theta_{108}c^5t_0^2/\hbar\) y
\(\Spec\mathcal G_5=\{G\text{ con multiplicidad }5,
0\text{ con multiplicidad }7\}\).
\textbf{Estatuto.} Selector: \emph{Teorema interno exacto}; realizaci\'on del portador:
\emph{Realizaci\'on f\'isica tipada}. \textbf{Criterio de refutación.} Otra ra\'iz positiva,
\(\Tr P_5\ne5\), o confundir cinco componentes con cinco polarizaciones
f\'isicas de GR sin masa.

\subsection*{Sistema de Planck, reciprocidad \(G\leftrightarrow\hbar\) e invariancia areal}
\textbf{Antecedente.} El producto bidireccional \(G_a\hbar_a\) y periodicidad temporal de (108) fases.
\textbf{Ecuaci\'on.}
\[
\ell_P^2=\frac{G\hbar}{c^3}=\theta_{108}\ell_0^2,
\quad t_P=\frac{54}{\pi}t_0,
\quad E_P=\frac{\pi}{54}\frac{\hbar}{t_0},
\]
\[
F_P=\frac{c^4}{G},\qquad P_P=\frac{c^5}{G},\qquad
\rho_P=\frac{\hbar}{\theta_{108}^2c^5t_0^4}.
\]
\textbf{Estructura y significado.} El \'area de Planck es invariante bajo el
desdoblamiento de acci\'on y enlaza gravedad con la escala t\'ermica.
\textbf{Valor.} Las secciones exactas anteriores.
\textbf{Estatuto.} \emph{Teorema interno exacto}. \textbf{Criterio de refutación.} Combinar
\(G\) y \(\hbar\) de hojas incompatibles o fallar la homogeneidad.

\subsection*{Operador de área y Hamiltoniano modular: \(\widehat{\mathcal A}_{\rm HMT}\), \(H_{\rm mod}\) y canales \(90/120\)}
\textbf{Antecedente.} Incidencia \(6\to8\), canales \(90/120\), salida
\(\gammaH\) y \(\ell_P^2\). \textbf{Ecuaci\'on.}
\[
\frac{A_{\rm phys}}{\ell_P^2}=\gammaH a_0(N_{90}+N_{120}),
\qquad
\HHM=\Dmod(90N_{90}+120N_{120})+cI.
\]
Si \(N=N_{90}+N_{120}\) y
\(D=90N_{90}+120N_{120}\),
\[
N_{90}=4N-D/30,\qquad N_{120}=D/30-3N.
\]
\textbf{Estructura y significado.} El \'area determina la ocupaci\'on total; el
Hamiltoniano conserva la procedencia de canal. Juntos reconstruyen el borde.
\textbf{Valor.} \(\gammaH a_0=0.0016755940749224991\ldots\) en la carta
declarada. \textbf{Estatuto.} \emph{Teorema interno exacto en corte finito}; no se
rederiva \(\gammaH\). \textbf{Criterio de refutación.} Conmutador no nulo, determinante
\(30\) perdido o uso del adjetivo hologr\'afico sin transporte al borde.

\subsection*{Entropía relativa, purificación termocampo doble y primera ley modular areal}
\textbf{Antecedente.} Estado reducido fiel del borde y su purificaci\'on de
Schmidt. \textbf{Ecuaci\'on.}
\[
I_{\rm or}=36\log3-S(\rho_A),
\quad
|\mathrm{TFD}\rangle=\sum_n\sqrt{p_n}|n\rangle|n\rangle,
\]
\[
\delta S=\beta_{90}\delta\langle \mathcal A_{90}\rangle
+\beta_{120}\delta\langle \mathcal A_{120}\rangle,
\qquad \frac{\beta_{120}}{\beta_{90}}=\frac43.
\]
\textbf{Estructura y significado.} La purificaci\'on convierte la entrop\'ia
reducida en entrelazamiento; la primera ley ecuaciona memoria modular y \'area.
\textbf{Valor.} \(I_{\rm or}=0.0016691832338689407\ldots\),
\(\beta_{90}=5.97264854010709\ldots\),
\(\beta_{120}=7.96353138680946\ldots\).
\textbf{Estatuto.} \emph{Teorema interno exacto local}. \textbf{Criterio de refutación.} Traza
parcial distinta de \(\rho_A\), \(I_{\rm or}<0\), o uso para variaciones
finitas sin el t\'ermino de entrop\'ia relativa.

\subsection*{L\'imite modular y par\'ametro de tipo III }
\textbf{Antecedente.} Estados locales compatibles con persistencia infinita de
ambos canales. \textbf{Ecuaci\'on.}
\[
\HHM^{(m)}=-\log\rho_m,
\qquad
\lambda_\partial=\exp(-30\Dmod).
\]
\textbf{Estructura y significado.} En el infinito vive una red local o el
generador modular GNS, no una matriz densidad global dentro de la UHF.
\textbf{Valor.}
\(\lambda_\partial=0.9966696464689573786\ldots\).
\textbf{Estatuto.} Red: \emph{Regla de tipo operatorio}; clasificaci\'on:
\emph{Teorema con hip\'otesis de persistencia}. \textbf{Criterio de refutación.} Ausencia
asint\'otica de un canal, estado no producto o raz\'on asint\'otica distinta.

\section{Involuci\'on de hojas, mezcla y sector electrod\'ebil}

\subsection*{Involución de hojas y transformación racional CKM: unitariedad, determinante y carácter de Jarlskog}
\textbf{Antecedente.} Intercambio \(C^*\mapsto-C^*\) y dato
\((h_6,o_8,t)=(6,8,3)\). \textbf{Ecuaci\'on.}
\[
R_{\rm CKM}=M\operatorname{diag}(1,-1,1)M^{-1},
\quad R_{\rm CKM}^2=I,\quad\det R_{\rm CKM}=-1,
\]
\[
-\frac{o_8-h_6}{t}=-\frac23,
\quad
A=\frac{1528\theta_{12}+3380\theta_{23}+801\theta_{13}}{3857},
\]
\[
C^*=6\frac{75\theta_{12}+176\theta_{23}-150\theta_{13}}{3857}.
\]
\textbf{Estructura y significado.} Barbero y CKM son realizaciones
funcionales distintas del cambio \(6\to8\), adoptado como antecedente
com\'un; la involuci\'on induce una reflexi\'on racional de rango impar uno.
\textbf{Valor.} Matriz racional exacta, \(\det M=-3857/6\); la fase CP
queda fija bajo esta involuci\'on, que no es conjugaci\'on CP.
\textbf{Estatuto.} \emph{Exacto geneal\'ogico/interno}.
\textbf{Criterio de refutación.} \(R^2\ne I\), fallo de \(RM=MD\), p\'erdida de
\(-2/3\), o presentar CKM como causa retrospectiva de \(\gammaH\).

\subsection*{Determinante angular, parámetro de Weinberg y acoplamientos electrodébiles}
\textbf{Antecedente.} Dos funciones del mismo operador angular:
\(F(T)^2\) y \(\det T\); rutas \(W/Z\) adyacentes.
\textbf{Ecuaci\'on.}
\[
D_A=\det T=e^{-2A_{\rm rad}},
\qquad \sin^2\theta_W^{\rm OS}=1-D_A,
\]
\[
g=2\sqrt{\frac{\pi\alpha}{1-D_A}},
\qquad g'=2\sqrt{\frac{\pi\alpha}{D_A}}.
\]
\textbf{Estructura y significado.} El vac\'io y Weinberg son realizaciones
espectrales distintas de un operador com\'un; el determinante es par bajo el
cambio de hoja.

\textbf{Valor.} \(D_A=0.7751291192471564302\ldots\) y
\(\sin^2\theta_W=0.2248708807528435698\ldots\).

\textbf{Estatuto.} Relaci\'on angular: \emph{Teorema interno exacto en la carta};
realizaci\'on de calibre: \emph{Realizaci\'on condicionada}.
\textbf{Criterio de refutación.} Rutas no adyacentes, cociente de masas
distinto de \(D_A\), o esquema de calibre no
declarado.

\subsection*{Constante de Fermi y acoplamiento de Higgs }
\textbf{Antecedente.} \(\alpha_{\HMT}\), \(D_A\), rutas \(W/H\) y carta
electrod\'ebil de \'arbol. \textbf{Ecuaci\'on.}
\[
G_F^{(0)}=\frac{\pi\alpha}{\sqrt2(1-D_A)m_W^2},
\]
\[
\lambda_H=\frac{\pi\alpha}{2(1-D_A)}
\exp\!\left(6A_{\rm rad}+\frac{10}{3}C^*_{\rm rad}\right),
\qquad G_F^{(0)}m_H^2=\sqrt2\lambda_H.
\]
\textbf{Estructura y significado.} Fermi es el l\'imite efectivo de baja
energ\'ia del canal cargado; Higgs conserva la orientaci\'on que el
determinante de calibre elimina.
\textbf{Valor.} \(G_F^{(0)}=1.1157162817\times10^{-5}\,\mathrm{GeV}^{-2}\)
en la carta declarada y \(\lambda_H=0.123847911548\ldots\).
\textbf{Estatuto.} \emph{Interfaz electrod\'ebil de \'arbol}.
\textbf{Criterio de refutación.} Usar \(G_F\) observado para elegir \(\Delta r\), omitir
esquema/escala o seleccionar \(\lambda_H\) desde la masa observada.

\subsection*{Anomal\'ia magn\'etica lept\'onica }
\textbf{Antecedente.} \((\alpha,\varphi,1000,54+1)\) con coeficientes
congelados antes del contraste. \textbf{Ecuaci\'on.}
\begin{equation}
a_\mu=\frac{\alpha}{2\pi}
+\frac{\alpha(\varphi-1)}{1000}
+\frac{\alpha^2}{1000(54+1)}.
\label{eq:atlas-anom-mu}
\end{equation}
\textbf{Estructura y significado.} Es una evaluaci\'on de tres funcionales HMT,
no una sustituci\'on de la serie radiativa completa de QED.
\textbf{Valor.} \(a_\mu=0.001165920713\ldots\).
\textbf{Estatuto.} \emph{Evaluaci\'on exacta con procedencia local abierta}.
\textbf{Criterio de refutación.} Elegir \(1000\) o \(55\) desde el valor experimental.

\section{Determinación de los parámetros cosmológicos mediante memoria torsional y dinámica de Perron}

\subsection*{Realización de de Sitter mediante la periodicidad temporal de \(108\) fases}
\textbf{Antecedente.} Selector circular \(R_{108}=\ell_P\) y realizaci\'on de
Sitter. \textbf{Ecuaci\'on.}
\[
H_{108}=\frac{\pi}{54t_0},\quad
\Lambda_{108}=\frac{\pi^2}{972\ell_0^2},\quad
\frac{T_{\rm dS}}{\Theta_{\rm clk}}=\frac{\log3}{2\pi},\quad
\frac{S_{\rm dS}}{k_B}=\pi.
\]
\textbf{Estructura y significado.} Un mismo radio determina expansi\'on,
curvatura, temperatura y entrop\'ia. \textbf{Valor.} El cu\'adruple exacto
anterior. \textbf{Estatuto.} \emph{Realizaci\'on f\'isica condicionada}.
\textbf{Criterio de refutación.} Fondo no de Sitter o realizador PCH que seleccione otro
radio.

\subsection*{Realizaci\'on de Perron--Einstein--Cartan }
\textbf{Antecedente.} Escala \(V_n=\varphi^n\), \(D=3\), modo estable
\(\varphi^{-2n}\). \textbf{Ecuaci\'on.}
\[
\frac{a_n}{a_0}=\varphi^{n/3},\qquad
M_n=\left(\frac{a_n}{a_0}\right)^{-6},\qquad
H_\varphi=\frac{\log\varphi}{3t_0}.
\]
\textbf{Estructura y significado.} El exponente \(-6\) es memoria activa y
coincide con la diluci\'on cuadr\'atica de esp\'in de Einstein--Cartan.
\textbf{Valor.} \(H_\varphi t_0=0.1604039416865344825\ldots\) y la ley exacta
\(M_na_n^6=a_0^6\). \textbf{Estatuto.} Escala:
\emph{Teorema interno exacto}; lectura de esp\'in: \emph{Realizaci\'on condicionada}.
\textbf{Criterio de refutación.} Exponente distinto, par\'ametro temporal no
aditivo o amplitud f\'isica
elegida desde un rebote.

\section{Separación tipada entre los radios de Bohr y Born--Infeld y las escalas generadas internamente}

\subsection*{Radio de Bohr como escala electrodinámica derivada de \(\alpha\)}
\textbf{Antecedente.} Acci\'on, vac\'io, \(\alpha\) y masa electr\'onica de su
antecedente. \textbf{Ecuaci\'on.}
\begin{equation}
a_0=\frac{\hbar}{m_ec\alpha}.
\label{eq:atlas-bohr}
\end{equation}
\textbf{Significado y valor.} Es una composici\'on
descendente, no un selector independiente de esta monograf\'ia.
\textbf{Estatuto.} \emph{Fuera de alcance}. \textbf{Criterio de refutación.} Recalcular
\(m_e\) aqu\'i o usar \(a_0\) para seleccionar \(\alpha\).

\subsection*{Escala de Born--Infeld y recuperación del límite de Maxwell}
\textbf{Antecedente.} \((h_{\rm int},c_{\rm int},e_{\rm int},\ell_0)\).
\textbf{Ecuaci\'on.}
\begin{equation}
E_*=\frac{h_{\rm int}c_{\rm int}}{e_{\rm int}\ell_0^2}.
\label{eq:atlas-born}
\end{equation}
\textbf{Estructura y significado.} La expresi\'on construye una escala de
campo; no construye todav\'ia una acci\'on Born--Infeld.
\textbf{Valor.} Secci\'on exacta \(E_*\), sin determinaci\'on de un
par\'ametro f\'isico.
\textbf{Estatuto.} \emph{Programa abierto}. \textbf{Criterio de refutación.} Identificar
por nombre o decimal sin acci\'on determinantal y selector.

\begin{proposition}[Resultado principal]
El resultado metrol\'ogico no es el decimal aislado. Es la conservaci\'on de
una genealog\'ia a trav\'es de operadores, rectas de magnitud y cartas, hasta
que el valor real aparece en la \'ultima posici\'on de la ficha.
\end{proposition}
